% GENERATED FILE. Edit canonical lessons, then run npm run generate:books.
% canonical-source-hash: fnv1a64:7829237a27ec603f
% canonical-lessons: KA-C396-masude, KA-C396-angikarisu, KA-C396-tiddupadi, KA-C396-sasaka, KA-C396-anumodisu

\chapter{Making Laws}
\label{ch:ka-making-laws}
\hlchaptermodality{396}

\begin{chapteropening}
\textbf{By the end of this chapter:} \emph{I can say a draft law, to approve, an amendment, a legislator and to endorse.}

Complete the last lesson of chapter 396: I can say a draft law, to approve, an amendment, a legislator and to endorse.
\end{chapteropening}

\section[\texorpdfstring{\kn{ಮಸೂದೆ} (masūde)}{ಮಸೂದೆ (masūde)}]{\kn{ಮಸೂದೆ} (masūde) — a draft law; a bill}

\label{lesson:KA-C396-masude}



\begin{quote}
Before the new one: say the Kannada for to implement, then the Kannada for a committee.
\end{quote}

\subsection*{You'll want to know: \kn{ಮಸೂದೆ}}
\textbf{\kn{ಮಸೂದೆ}} — \emph{masūde} — \textquotedblleft{}a draft law; a bill\textquotedblright{}.

\begin{quote}
\kn{ಮಂತ್ರಿ} \kn{ಹೊಸ} \kn{ಮಸೂದೆಯ} \kn{ಬಗ್ಗೆ} \kn{ಮಾತಾಡಿದರು}.
\end{quote}

\begin{quote}
\emph{mantri hosa masūdeya bagge mātāḍidaru.}
\end{quote}

\begin{quote}
The minister spoke about the new bill.
\end{quote}

A \kn{ಮಸೂದೆ} becomes a \kn{ಕಾಯ್ದೆ} once it is passed.

\begin{cousinweb}
\textbf{\kn{ಮಸೂದೆ}} (\emph{masūde}) is usually traced to Arabic \emph{musawwada}, \textquotedblleft{}a draft\textquotedblright{}, through Hindustani \emph{masauda}.
\end{cousinweb}

\begin{grammarlens}[title={What you've built}]
One more word for giving reasons and opinions.
\end{grammarlens}

\subsection*{Your turn}
\begin{itemize}
\raggedright
  \item \emph{Say it:} \emph{masūde}
  \item \emph{Say it:} \emph{masūde}, once more
  \item \emph{Say it:} \emph{mantri hosa masūdeya bagge mātāḍidaru.} — The minister spoke about the new bill.
  \item \emph{Recall:} say the Kannada for to implement, then the Kannada for a committee, then say \emph{masūde} again
\end{itemize}

\begin{tcolorbox}[breakable,colback=teal!4,colframe=teal!35!black,title={Before you move on}]
What does \kn{ಮಸೂದೆ} mean? (\textquotedblleft{}A draft law; a bill\textquotedblright{}.) Say it once more.
\end{tcolorbox}
\section[\texorpdfstring{\kn{ಅಂಗೀಕರಿಸು} (aṅgīkarisu)}{ಅಂಗೀಕರಿಸು (aṅgīkarisu)}]{\kn{ಅಂಗೀಕರಿಸು} (aṅgīkarisu) — to approve; to pass}

\label{lesson:KA-C396-angikarisu}



\begin{quote}
Before the new one: say the Kannada for a committee, then the Kannada for a draft law.
\end{quote}

\subsection*{You'll want to know: \kn{ಅಂಗೀಕರಿಸು}}
\textbf{\kn{ಅಂಗೀಕರಿಸು}} — \emph{aṅgīkarisu} — \textquotedblleft{}to approve; to pass\textquotedblright{}.

\begin{quote}
\kn{ಸಂಸತ್ತು} \kn{ಮಸೂದೆಯನ್ನು} \kn{ಅಂಗೀಕರಿಸಿತು}.
\end{quote}

\begin{quote}
\emph{saṃsattu masūdeyannu aṅgīkarisitu.}
\end{quote}

\begin{quote}
Parliament passed the bill.
\end{quote}

Past (I …): \kn{ಅಂಗೀಕರಿಸಿದೆ}, \emph{aṅgīkariside}.

\kn{ಅಂಗೀಕಾರ}, approval.

\begin{cousinweb}
\textbf{\kn{ಅಂಗೀಕರಿಸು}} (\emph{aṅgīkarisu}) is built on Sanskrit \emph{aṅgīkāra}, \textquotedblleft{}agreement, promise\textquotedblright{} (Monier-Williams), with the Kannada verb ending -\emph{isu}.
\end{cousinweb}

\begin{grammarlens}[title={What you've built}]
One more word for giving reasons and opinions.
\end{grammarlens}

\subsection*{Your turn}
\begin{itemize}
\raggedright
  \item \emph{Say it:} \emph{aṅgīkarisu}
  \item \emph{Say it:} \emph{aṅgīkarisu}, once more
  \item \emph{Say it:} \emph{saṃsattu masūdeyannu aṅgīkarisitu.} — Parliament passed the bill.
  \item \emph{Recall:} say the Kannada for a committee, then the Kannada for a draft law, then say \emph{aṅgīkarisu} again
\end{itemize}

\begin{tcolorbox}[breakable,colback=teal!4,colframe=teal!35!black,title={Before you move on}]
What does \kn{ಅಂಗೀಕರಿಸು} mean? (\textquotedblleft{}To approve; to pass\textquotedblright{}.) Say it once more.
\end{tcolorbox}
\section[\texorpdfstring{\kn{ತಿದ್ದುಪಡಿ} (tiddupaḍi)}{ತಿದ್ದುಪಡಿ (tiddupaḍi)}]{\kn{ತಿದ್ದುಪಡಿ} (tiddupaḍi) — an amendment}

\label{lesson:KA-C396-tiddupadi}



\begin{quote}
Before the new one: say the Kannada for a draft law, then the Kannada for to approve.
\end{quote}

\subsection*{You'll want to know: \kn{ತಿದ್ದುಪಡಿ}}
\textbf{\kn{ತಿದ್ದುಪಡಿ}} — \emph{tiddupaḍi} — \textquotedblleft{}an amendment\textquotedblright{}.

\begin{quote}
\kn{ಕಾನೂನಿಗೆ} \kn{ಒಂದು} \kn{ತಿದ್ದುಪಡಿ} \kn{ಬೇಕು}.
\end{quote}

\begin{quote}
\emph{kānūnige ondu tiddupaḍi bēku.}
\end{quote}

\begin{quote}
The law needs an amendment.
\end{quote}

From \kn{ತಿದ್ದು}, to correct. \kn{ತಿದ್ದುಪಡಿ} \kn{ತರು}, to bring an amendment.

\begin{cousinweb}
\textbf{\kn{ತಿದ್ದುಪಡಿ}} (\emph{tiddupaḍi}) is native Dravidian, built on \emph{tiddu}, \textquotedblleft{}to make straight, to correct\textquotedblright{} (DEDR 3251).
\end{cousinweb}

\begin{grammarlens}[title={What you've built}]
One more word for giving reasons and opinions.
\end{grammarlens}

\subsection*{Your turn}
\begin{itemize}
\raggedright
  \item \emph{Say it:} \emph{tiddupaḍi}
  \item \emph{Say it:} \emph{tiddupaḍi}, once more
  \item \emph{Say it:} \emph{kānūnige ondu tiddupaḍi bēku.} — The law needs an amendment.
  \item \emph{Recall:} say the Kannada for a draft law, then the Kannada for to approve, then say \emph{tiddupaḍi} again
\end{itemize}

\begin{tcolorbox}[breakable,colback=teal!4,colframe=teal!35!black,title={Before you move on}]
What does \kn{ತಿದ್ದುಪಡಿ} mean? (\textquotedblleft{}An amendment\textquotedblright{}.) Say it once more.
\end{tcolorbox}
\section[\texorpdfstring{\kn{ಶಾಸಕ} (śāsaka)}{ಶಾಸಕ (śāsaka)}]{\kn{ಶಾಸಕ} (śāsaka) — a legislator}

\label{lesson:KA-C396-sasaka}



\begin{quote}
Before the new one: say the Kannada for to approve, then the Kannada for an amendment.
\end{quote}

\subsection*{You'll want to know: \kn{ಶಾಸಕ}}
\textbf{\kn{ಶಾಸಕ}} — \emph{śāsaka} — \textquotedblleft{}a legislator\textquotedblright{}.

\begin{quote}
\kn{ನಮ್ಮ} \kn{ಶಾಸಕರು} \kn{ಊರಿನ} \kn{ರಸ್ತೆ} \kn{ಸರಿಪಡಿಸಿದರು}.
\end{quote}

\begin{quote}
\emph{namma śāsakaru ūrina raste saripaḍisidaru.}
\end{quote}

\begin{quote}
Our legislator got the town road repaired.
\end{quote}

Plural: \kn{ಶಾಸಕರು}, \emph{śāsakaru}.

A member of a state assembly. \kn{ಶಾಸನ}, a law or an inscription.

\begin{cousinweb}
\textbf{\kn{ಶಾಸಕ}} (\emph{śāsaka}) is borrowed from Sanskrit \emph{śāsaka}, \textquotedblleft{}a ruler, a governor, an instructor\textquotedblright{} (Monier-Williams).
\end{cousinweb}

\begin{grammarlens}[title={What you've built}]
One more word for giving reasons and opinions.
\end{grammarlens}

\subsection*{Your turn}
\begin{itemize}
\raggedright
  \item \emph{Say it:} \emph{śāsaka}
  \item \emph{Say it:} \emph{śāsaka}, once more
  \item \emph{Say it:} \emph{namma śāsakaru ūrina raste saripaḍisidaru.} — Our legislator got the town road repaired.
  \item \emph{Recall:} say the Kannada for to approve, then the Kannada for an amendment, then say \emph{śāsaka} again
\end{itemize}

\begin{tcolorbox}[breakable,colback=teal!4,colframe=teal!35!black,title={Before you move on}]
What does \kn{ಶಾಸಕ} mean? (\textquotedblleft{}A legislator\textquotedblright{}.) Say it once more.
\end{tcolorbox}
\section[\texorpdfstring{\kn{ಅನುಮೋದಿಸು} (anumōdisu)}{ಅನುಮೋದಿಸು (anumōdisu)}]{\kn{ಅನುಮೋದಿಸು} (anumōdisu) — to endorse}

\label{lesson:KA-C396-anumodisu}



\begin{quote}
Before the new one: say the Kannada for an amendment, then the Kannada for a legislator.
\end{quote}

\subsection*{You'll want to know: \kn{ಅನುಮೋದಿಸು}}
\textbf{\kn{ಅನುಮೋದಿಸು}} — \emph{anumōdisu} — \textquotedblleft{}to endorse\textquotedblright{}.

\begin{quote}
\kn{ಸಭೆಯಲ್ಲಿ} \kn{ಎಲ್ಲರೂ} \kn{ಅವಳ} \kn{ಸಲಹೆಯನ್ನು} \kn{ಅನುಮೋದಿಸಿದರು}.
\end{quote}

\begin{quote}
\emph{sabheyalli ellarū avaḷa salaheyannu anumōdisidaru.}
\end{quote}

\begin{quote}
At the meeting everyone endorsed her suggestion.
\end{quote}

Past (I …): \kn{ಅನುಮೋದಿಸಿದೆ}, \emph{anumōdiside}.

\kn{ಅನುಮೋದನೆ}, approval.

\begin{cousinweb}
\textbf{\kn{ಅನುಮೋದಿಸು}} (\emph{anumōdisu}) is built on Sanskrit \emph{anumodana}, \textquotedblleft{}applauding, assent\textquotedblright{} (Monier-Williams), with the Kannada verb ending -\emph{isu}.
\end{cousinweb}

\begin{grammarlens}[title={What you've built}]
One more word for giving reasons and opinions.
\end{grammarlens}

\subsection*{Your turn}
\begin{itemize}
\raggedright
  \item \emph{Say it:} \emph{anumōdisu}
  \item \emph{Say it:} \emph{anumōdisu}, once more
  \item \emph{Say it:} \emph{sabheyalli ellarū avaḷa salaheyannu anumōdisidaru.} — At the meeting everyone endorsed her suggestion.
  \item \emph{Recall:} say the Kannada for an amendment, then the Kannada for a legislator, then say \emph{anumōdisu} again
\end{itemize}

\begin{tcolorbox}[breakable,colback=teal!4,colframe=teal!35!black,title={Before you move on}]
What does \kn{ಅನುಮೋದಿಸು} mean? (\textquotedblleft{}To endorse\textquotedblright{}.) Say it once more.
\end{tcolorbox}
